\section{Test działania wykrywacza}

Podczas pracy wykrywacza pobierany prąd zawarty był w zakresie $(70mA - 80mA)$, przy przyjętym wcześniej napięciu zasilania $12V$ daje to stałe obciążenie około $0,9W$.

Urządzenie poprawnie wyznacza czas trwania odpowiedzi sondy, wraz z filtracją szumów. Na rysunku \ref{fig:filtry_wewnetrzne} wykreślono surowe, pomierzone dane próbek wraz ze stanami filtrów EMA oraz końcowym, wynikowym sygnałem
będący podstawą do podjęcia decyzji o stwierdzonej anomalii sygnału. Na wykresie \ref{fig:filtry_wykryte} zauważyć można zmianę wartości spowodowaną obecnością obiektu badanego w pobliżu.


\begin{figure}
    \centering
    \includesvg[width=\textwidth]{img/testy/filtry.svg}
    \caption{Wykres wewnętrznych stanów wykrywacza}\label{fig:filtry_wewnetrzne}
\end{figure}

\begin{figure}
    \centering
    \includesvg[width=\textwidth]{img/testy/filtry_wykryte.svg}
    \caption{Wykres wewnętrznych stanów wykrywacza ze zbliżonym tymczasowo obiektem}\label{fig:filtry_wykryte}
\end{figure}


